\section{Discussion and Conclusion}
\label{sec:discussion}

Perturbation magnitude tells us how strongly a model responds, but responses
of the same size can reflect different behaviors. \decaf uses the final
factual--counterfactual contrast to separate intermediate responses into
evidence, contradiction, and fragility, while preserving ordinary magnitude
exactly, $\Abs=E+C+F$.

The experiments show that these distinctions are not artifacts of the
definitions. They track independently measured behavior in controlled vision
and tabular settings, remain distinguishable on natural images after magnitude
is controlled, and remain useful under external attribution criteria.
ImageNet further shows why this matters: changing only the reveal path can
greatly increase total response without increasing evidence.

These response roles are relative to the chosen counterfactual pair, score,
and reveal path; they are not claims about hidden causal mechanisms. The
endpoint most directly describes its own intervention, while the trajectory is
most useful beyond that intervention. Magnitude tells us how much a model
reacts; \decaf preserves that quantity while revealing what kind of response
produced it.